\documentclass{article}
\usepackage[margin=0.8in]{geometry}
\usepackage{titling}

% BIBLIOGRAPHY
\usepackage[
backend=bibtex,
style=numeric,
]{biblatex}
\addbibresource{references.bib}

% TEXT, EQUATIONS, AND TABLES
\usepackage[none]{hyphenat}
\usepackage{amsmath}
\usepackage{booktabs}
\usepackage{float}
\usepackage{hyperref}
\usepackage{microtype}
\numberwithin{equation}{section}
\setlength{\parindent}{0pt}
\sloppy

% FIGURES
\usepackage{graphicx, wrapfig}
\usepackage{subcaption}
\usepackage{caption}
\graphicspath{{./plots/}{./images/}}
\captionsetup[figure]{font=footnotesize,labelfont=footnotesize}

\begin{document}

\begin{center}
{\huge\bfseries Statistical Analysis of the Exoplanet Mass--Radius Relation\\}
\vspace{0.5cm}
{\Large ETH Z\"urich, Switzerland\\}
\vspace{0.25cm}
Team Project (supervisor: Dr. T. Tr\"oster)\\
September 2024--January 2025\\
\vspace{0.25cm}
Guglielmo De Toma
\end{center}
\vspace{5pt}

\begin{abstract}
\noindent
This project compares Bayesian and machine-learning descriptions of the exoplanet mass--radius ($M$--$R$) relation using the DACE catalogue \cite{dace_catalog}. The analysis fits $\log_{10}(R/R_\oplus)$ as a function of $\log_{10}(M/M_\oplus)$. A continuous broken power law explicitly includes radius uncertainty, propagated mass uncertainty, and segment-dependent intrinsic scatter. The baseline two-breakpoint \texttt{emcee} fit gives slopes $b_1=0.397$, $b_2=0.613$, and $b_3=-0.035$, identifying small-planet, intermediate, and giant-planet regimes. Its posterior predictive check is well calibrated, with PTE = 0.524 and 0.963 of planets within two total standard deviations. BIC favours two breakpoints, WAIC favours more flexible three- and four-breakpoint models, and the current \texttt{dynesty} evidence calculation favours two breakpoints. A Rational Quadratic Gaussian Process selected by training-only cross-validation gives held-out RMSE = 0.1069 dex, close to the two-breakpoint Bayesian model's 0.1063 dex on the identical test set. The selected Random Forest gives RMSE = 0.1192 dex. Both flexible methods have imperfect predictive calibration, especially the Random Forest, whose tree-spread uncertainty is only an ensemble-disagreement proxy. The shared comparison therefore finds that the broken-power-law model offers the best combination of physical interpretation and probabilistic performance, while the GP is a competitive non-parametric predictor.
\end{abstract}

\section{Introduction}

Since the discovery of the first exoplanet orbiting a Sun-like star in 1995 \cite{mayor1995jupiter}, exoplanet science has rapidly evolved from detection toward physical characterization. More than 5500 exoplanets have now been discovered, revealing a broad diversity of planetary systems and planetary properties. Among the quantities used to characterize exoplanets, planetary mass and radius play a central role because together they constrain bulk density, composition, internal structure, and formation history.

Planetary masses and radii are usually measured through different observational techniques, and many exoplanets have only one of the two quantities available. The mass--radius relation is therefore a crucial statistical tool: it can be used to infer missing planetary properties, classify planets into broad physical populations, and investigate the physics governing planetary interiors and atmospheres.

The exoplanet $M$--$R$ relation is non-linear and intrinsically scattered. Terrestrial planets, Neptune-like planets, and gas giants do not follow a single power law because their radii respond differently to changes in mass. Moreover, observational uncertainty, composition, atmospheric mass fraction, irradiation, age, and evolutionary history all contribute to the observed dispersion. Modelling the relation is therefore both a physical and a statistical problem.

In this work, the relation is analysed with complementary Bayesian and machine-learning approaches: broken-power-law inference sampled with \texttt{emcee}, Bayesian evidence estimation with \texttt{dynesty}, Gaussian Process (GP) regression, and Random Forest (RF) regression. A final leakage-free pipeline evaluates the MCMC, GP, and RF predictions on exactly the same held-out planets. The goal is to compare fit quality, model selection, uncertainty quantification, physical interpretability, and extrapolation behaviour without conflating in-sample fit with predictive performance.

\section{Bayesian Analysis with \texttt{emcee}}

\subsection{Dataset and Preprocessing}

The Bayesian analysis was performed using the DACE exoplanet catalogue \cite{dace_catalog}. Only planets with measured masses, radii, and associated uncertainties were retained. Planetary masses and radii were converted from Jupiter units to Earth units using

\begin{equation}
M = M_{\mathrm{Jup}} \times 317.8284,
\end{equation}

\begin{equation}
R = R_{\mathrm{Jup}} \times 11.2089.
\end{equation}

The relation was fitted in logarithmic space,

\begin{equation}
x = \log_{10}\left(\frac{M}{M_\oplus}\right),
\qquad
y = \log_{10}\left(\frac{R}{R_\oplus}\right).
\end{equation}

This transformation is useful because a power law,

\begin{equation}
R \propto M^\alpha,
\end{equation}

becomes a linear relation in log-log space. The slope of the line is therefore directly interpretable as the power-law exponent. Measurement uncertainties were propagated into logarithmic space with first-order error propagation:

\begin{equation}
\sigma_{\log M} =
\frac{\sigma_M}{M \ln 10},
\qquad
\sigma_{\log R} =
\frac{\sigma_R}{R \ln 10}.
\end{equation}

No radius or mass cut was applied. This is important because the purpose of the broken-power-law model is precisely to capture the different regimes present in the full exoplanet population.

\subsection{Piecewise-Linear Model}

The model assumes that $\log_{10} R$ is a continuous piecewise-linear function of $\log_{10} M$. For the baseline case with two breakpoints,

\begin{equation}
f(x)
=
a + b_1 x
+
(b_2-b_1)\max(0,x-x_1)
+
(b_3-b_2)\max(0,x-x_2),
\end{equation}

where $a$ is the intercept, $b_i$ are the slopes of the three regimes, and $x_i$ are the breakpoint locations in $\log_{10}(M/M_\oplus)$. The hinge-function formulation enforces continuity while allowing the slope to change between regimes. Because the model is fitted in log-log space, each slope corresponds to

\begin{equation}
R \propto M^{b_i}.
\end{equation}

The same implementation was also run with one, three, and four breakpoints in order to test whether additional flexibility was justified by the data.

\subsection{Intrinsic Scatter and Likelihood}

The observed spread in radius at fixed mass is larger than expected from measurement errors alone. To account for this, each segment has its own intrinsic scatter term. For planet $i$, the total variance in the $y$ direction is

\begin{equation}
\sigma_{\mathrm{tot},i}^2
=
\sigma_{y,i}^2
+
\left[
\frac{dy}{dx}(x_i)\sigma_{x,i}
\right]^2
+
\sigma_{\mathrm{int},s(i)}^2,
\end{equation}

where $s(i)$ denotes the segment containing the planet. The second term propagates the mass uncertainty into the radius direction using the local slope of the fitted relation. The intrinsic scatter parameters are sampled in logarithmic form,

\begin{equation}
\sigma_{\mathrm{int},s} = \exp(\log \sigma_{\mathrm{int},s}),
\end{equation}

which guarantees positivity.

The likelihood is Gaussian in log-radius:

\begin{equation}
y_i \sim \mathcal{N}\left(f(x_i;\theta),\sigma_{\mathrm{tot},i}^2\right),
\end{equation}

with log-likelihood

\begin{equation}
\log \mathcal{L}
=
-\frac{1}{2}
\sum_i
\left[
\frac{(y_i-f(x_i;\theta))^2}{\sigma_{\mathrm{tot},i}^2}
+
\log(2\pi \sigma_{\mathrm{tot},i}^2)
\right].
\end{equation}

The logarithmic variance term is essential: it penalizes arbitrarily large scatter values, so the likelihood cannot be improved simply by inflating $\sigma_{\mathrm{int}}$.

\subsection{Priors and Sampling}

Uniform priors were adopted within broad bounds for the intercept, slopes, breakpoints, and intrinsic scatters. The breakpoints were constrained to remain ordered and separated by a minimum distance in $\log_{10}M$, avoiding degenerate solutions in which two segments collapse onto each other.

Before posterior sampling, a maximum a posteriori estimate was obtained with the L-BFGS-B optimizer using multiple deterministic and randomized starting points. The MCMC chains were then initialized around this solution and sampled with the affine-invariant ensemble sampler implemented in \texttt{emcee} \cite{foremanmackey2013emcee}. The posterior distribution is

\begin{equation}
p(\theta|D) \propto p(D|\theta)p(\theta).
\end{equation}

After sampling, burn-in and thinning were estimated from the autocorrelation time when possible.

\subsection{Posterior Predictive Checks and Model Comparison}

Posterior predictive checks were used to test whether the fitted model can reproduce the observed scatter. For each posterior draw, replicated data were generated as

\begin{equation}
y_i^{\mathrm{rep}}
\sim
\mathcal{N}\left(f(x_i;\theta),\sigma_{\mathrm{tot},i}^2\right).
\end{equation}

The test statistic was

\begin{equation}
\chi^2
=
\sum_i
\frac{(y_i-f(x_i;\theta))^2}{\sigma_{\mathrm{tot},i}^2}.
\end{equation}

The posterior predictive $p$-value, or probability-to-exceed (PTE), was computed as the fraction of replicated datasets with a larger test statistic than the observed data. Values near 0 or 1 would indicate tension; values closer to the middle of the unit interval indicate that the observed data are typical under the fitted model.

Models with different numbers of breakpoints were compared using BIC and WAIC. The BIC is

\begin{equation}
\mathrm{BIC}
=
k\ln N
-
2\log \mathcal{L}_{\mathrm{max}},
\end{equation}

where $k$ is the number of parameters and $N$ is the number of planets. WAIC was estimated from pointwise log-likelihood values over posterior samples and is used as an estimate of out-of-sample predictive performance.

\subsection{Results}

The baseline two-breakpoint fit used 760 planets. The posterior median breakpoints were
$x_1 = 0.824^{+0.221}_{-0.006}$ and
$x_2 = 2.181^{+0.020}_{-0.010}$, corresponding to approximately
$6.7\,M_\oplus$ and $152\,M_\oplus$. These values separate a small-planet regime, an intermediate regime, and a giant-planet regime.

\begin{table}[H]
\centering
\caption{Posterior summary for the baseline two-breakpoint \texttt{emcee} model. Slopes are power-law exponents in $R\propto M^b$. Intrinsic scatters are in dex.}
\label{tab:emcee-parameters}
\begin{tabular}{lccc}
\toprule
Quantity & Median & $-1\sigma$ & $+1\sigma$ \\
\midrule
$b_1$ & 0.397 & 0.031 & 0.034 \\
$b_2$ & 0.613 & 0.012 & 0.014 \\
$b_3$ & -0.035 & 0.010 & 0.010 \\
$x_1$ & 0.824 & 0.006 & 0.221 \\
$x_2$ & 2.181 & 0.010 & 0.020 \\
$\sigma_{\mathrm{int},1}$ & 0.071 & 0.010 & 0.011 \\
$\sigma_{\mathrm{int},2}$ & 0.122 & 0.006 & 0.008 \\
$\sigma_{\mathrm{int},3}$ & 0.072 & 0.003 & 0.003 \\
\bottomrule
\end{tabular}
\end{table}

Figure~\ref{fig:emcee-fit} shows the fitted relation. The blue bands represent uncertainty in the mean relation, while the green bands show the posterior predictive spread including intrinsic scatter. The coloured vertical bars mark the central 68\% and 95\% posterior intervals for the breakpoint locations, with dashed coloured lines at the posterior medians. This distinction is important: many observed planets are expected to lie outside the narrow mean-relation credible band, but they should be better described by the wider predictive band.

\begin{figure}[H]
\centering
\includegraphics[width=0.92\linewidth]{emcee_mass_radius_fit}
\caption{Bayesian fit of $\log_{10}(R/R_\oplus)$ as a function of $\log_{10}(M/M_\oplus)$ for the two-breakpoint model. The posterior median and MAP relation are shown together with credible bands for the mean relation, predictive bands including intrinsic scatter, and coloured 68\% and 95\% breakpoint intervals.}
\label{fig:emcee-fit}
\end{figure}

The three slopes have clear physical interpretations. The first segment has a positive exponent compatible with rocky and super-Earth planets. The intermediate segment is steeper, reflecting the rapid increase in radius associated with volatile-rich planets. The final slope is close to zero, consistent with the weak mass dependence of gas-giant radii over a broad mass range.

\begin{figure}[H]
\centering
\begin{subfigure}{0.48\linewidth}
\centering
\includegraphics[width=\linewidth]{emcee_mass_radius_trace}
\caption{Trace plot.}
\end{subfigure}
\hfill
\begin{subfigure}{0.48\linewidth}
\centering
\includegraphics[width=\linewidth]{emcee_mass_radius_corner}
\caption{Corner plot.}
\end{subfigure}
\caption{MCMC diagnostics for the baseline \texttt{emcee} model. The trace plot checks chain stability, while the corner plot shows posterior correlations between slopes, breakpoints, and intrinsic scatters.}
\label{fig:emcee-diagnostics}
\end{figure}

The posterior predictive check gives a PTE of 0.524 for the final two-breakpoint run. In addition, the MAP residual coverage test finds that 525 out of 760 planets, or 0.691 of the sample, lie within one total standard deviation of the fitted line, and 732 out of 760 planets, or 0.963, lie within two total standard deviations. These numbers are close to the Gaussian expectations of about 0.68 and 0.95, indicating that the combination of measurement errors, propagated mass errors, and intrinsic scatter is reasonably calibrated.

\begin{figure}[H]
\centering
\includegraphics[width=0.78\linewidth]{emcee_mass_radius_ppd_pte}
\caption{Posterior predictive check for the baseline two-breakpoint model. The histogram shows the difference between replicated and observed chi-square statistics; the vertical line marks zero. The corresponding PTE is 0.524.}
\label{fig:emcee-ppd}
\end{figure}

\begin{table}[H]
\centering
\caption{Model comparison for \texttt{emcee} fits with different numbers of breakpoints. Lower BIC and WAIC are preferred.}
\label{tab:emcee-model-comparison}
\begin{tabular}{ccccccc}
\toprule
Breakpoints & BIC & $\Delta$BIC & WAIC & $\Delta$WAIC & PTE & $f_{2\sigma}$ \\
\midrule
1 & -1419.19 & 35.10 & -1447.13 & 65.76 & 0.557 & 0.961 \\
2 & -1454.29 & 0.00 & -1494.35 & 18.54 & 0.553 & 0.958 \\
3 & -1453.63 & 0.66 & -1511.90 & 0.99 & 0.577 & 0.954 \\
4 & -1440.69 & 13.60 & -1512.89 & 0.00 & 0.427 & 0.963 \\
\bottomrule
\end{tabular}
\end{table}

The model-comparison statistics do not point to a single conclusion in a completely model-independent way. BIC strongly penalizes additional parameters and favours the two-breakpoint model, with the three-breakpoint model very close behind. WAIC instead prefers the four-breakpoint model, with the three-breakpoint model nearly indistinguishable. The posterior predictive PTE values are acceptable for all models. Since the two-breakpoint model captures the intended three broad planetary regimes while remaining parsimonious and physically interpretable, it is adopted as the baseline Bayesian model for comparison with the Gaussian Process and Random Forest approaches.

\subsection{Discussion}

The \texttt{emcee} analysis provides a physically interpretable probabilistic description of the mass--radius relation. Its main strength is that the model parameters map directly onto planetary-regime slopes, breakpoints, and intrinsic scatters. The likelihood also treats radius errors, propagated mass errors, and astrophysical scatter in a single coherent variance model.

The main limitation is that the functional form must be chosen in advance. A piecewise-linear relation is useful for identifying broad regimes, but it cannot capture every localized feature of the data. Increasing the number of breakpoints improves flexibility, as reflected by WAIC, but at the cost of interpretability and higher posterior dimension. This motivates the comparison with Gaussian Processes and Random Forests, which are more flexible but less directly tied to simple physical scaling laws.

\section{Nested Sampling with \texttt{dynesty}}

\subsection{Motivation}

The \texttt{emcee} analysis samples the posterior distribution for a chosen model, while BIC and WAIC are then used as approximate model-comparison statistics. Nested sampling provides a complementary Bayesian approach because it estimates the marginal likelihood, or evidence, directly. For model $\mathcal{M}$, the evidence is

\begin{equation}
Z_{\mathcal{M}}
=
p(D|\mathcal{M})
=
\int
\mathcal{L}(\theta)\pi(\theta)\,d\theta,
\end{equation}

where $\mathcal{L}(\theta)$ is the likelihood and $\pi(\theta)$ is the prior. The evidence rewards models that fit the data well over a substantial region of prior volume and penalizes models whose good fits occupy only a very small part of parameter space. Two models can therefore be compared through the Bayes factor,

\begin{equation}
\log B_{12}
=
\log Z_1 - \log Z_2.
\end{equation}

This is particularly useful for the mass--radius problem because changing the number of breakpoints changes both the goodness of fit and the effective prior volume.

\subsection{Implementation}

The nested-sampling analysis used the same cleaned DACE dataset, the same log-radius model, and the same total variance model as the \texttt{emcee} analysis. The likelihood remained

\begin{equation}
\log \mathcal{L}
=
-\frac{1}{2}
\sum_i
\left[
\frac{(y_i-f(x_i;\theta))^2}{\sigma_{\mathrm{tot},i}^2}
+
\log(2\pi \sigma_{\mathrm{tot},i}^2)
\right].
\end{equation}

The calculation was performed with \texttt{dynesty} \cite{speagle2020dynesty}. The unit-cube prior transform maps the intercept, slopes, and logarithmic intrinsic scatters to the same broad prior ranges used in the MCMC analysis. Breakpoints are drawn as ordered positions between the 5th and 95th percentiles of the observed log-mass distribution, with a minimum spacing imposed to avoid exactly degenerate segments.

The evidence comparison shown here used the broad prior, 500 live points, the \texttt{rwalk} sampling method, the \texttt{multi} bounding method, and a target stopping criterion of $\Delta\log Z = 0.1$. These settings are sufficient for a stable project-level comparison, although a final publication-quality evidence calculation should still check sensitivity to sampler settings and random seeds.

\subsection{Evidence Comparison}

Table~\ref{tab:dynesty-evidence} summarizes the nested-sampling evidence comparison. In this run, the two-breakpoint model has the largest evidence. The three-breakpoint and four-breakpoint models are lower by $\Delta\log Z=-5.26$ and $-6.89$, respectively, while the one-breakpoint model is strongly disfavoured. Because nested sampling is stochastic and the flexible models occupy larger prior volumes, this evidence ranking should be interpreted together with the posterior fit shapes rather than as a purely mechanical model choice.

\begin{table}[H]
\centering
\caption{Nested-sampling model comparison with \texttt{dynesty}. The table uses the broad prior. $\Delta\log Z$ is measured relative to the highest-evidence model.}
\label{tab:dynesty-evidence}
\begin{tabular}{cccccc}
\toprule
Breakpoints & $\log Z$ & $\sigma_{\log Z}$ & $\Delta\log Z$ & PTE & $f_{2\sigma}$ \\
\midrule
1 & 698.85 & 0.28 & -21.89 & 0.551 & 0.962 \\
2 & 720.75 & 0.31 & 0.00 & 0.537 & 0.963 \\
3 & 715.49 & 0.36 & -5.26 & 0.673 & 0.959 \\
4 & 713.86 & 0.36 & -6.89 & 0.891 & 0.964 \\
\bottomrule
\end{tabular}
\end{table}

The evidence result should be interpreted together with the shape of the posterior fit. Figure~\ref{fig:dynesty-fits} compares the three- and four-breakpoint nested-sampling fits, including coloured vertical bars for the central 68\% and 95\% posterior intervals of each breakpoint and dashed coloured lines for the posterior medians. The two fits are qualitatively similar, but the four-breakpoint model introduces an additional transition region whose evidence support is weak relative to the three-breakpoint model in this run.

\begin{figure}[H]
\centering
\begin{subfigure}{0.49\linewidth}
\centering
\includegraphics[width=\linewidth]{dynesty_mass_radius_3_breakpoints_broad_fit}
\caption{Three breakpoints.}
\end{subfigure}
\hfill
\begin{subfigure}{0.49\linewidth}
\centering
\includegraphics[width=\linewidth]{dynesty_mass_radius_4_breakpoints_broad_fit}
\caption{Four breakpoints.}
\end{subfigure}
\caption{Nested-sampling fits for the three- and four-breakpoint flexible models. Bands are defined as in the \texttt{emcee} fit: blue bands describe the posterior uncertainty in the mean relation, green bands describe the intrinsic posterior predictive spread, and coloured vertical bars show the central 68\% and 95\% breakpoint intervals.}
\label{fig:dynesty-fits}
\end{figure}

\subsection{Discussion}

Nested sampling and the MCMC information criteria answer related but not identical questions. The \texttt{emcee} comparison showed that BIC favours two breakpoints because BIC imposes a strong penalty on additional parameters. WAIC, which is more focused on predictive accuracy, preferred the more flexible three- and four-breakpoint models. The evidence from \texttt{dynesty} is closer in spirit to the full Bayesian model comparison: it integrates over the prior volume and therefore automatically includes an Occam penalty.

The nested-sampling result favours more than one regime, but in this run the evidence maximum returns to the parsimonious two-breakpoint model. The three- and four-breakpoint fits remain useful as flexible descriptions of possible additional structure, and their relative ordering can move with the stochastic nested-sampling run. Therefore, the main conclusion is not that any single breakpoint count is uniquely required. Rather, the data clearly reject a single-break model, support multiple planetary regimes, and leave some ambiguity between a parsimonious physically interpretable model and a slightly more flexible statistical description.

\section{Gaussian Process Regression}

\subsection{Motivation}

The broken-power-law models impose a specific functional form on the mass--radius relation. This makes the inferred slopes and breakpoints physically interpretable, but it also restricts the shape of the relation. Gaussian Process (GP) regression provides a complementary non-parametric approach: instead of choosing a fixed number of regimes, the function itself is treated as a random draw from a distribution over smooth functions \cite{rasmussen2006gaussian}.

In log-log space, the model can be written as

\begin{equation}
y(x) = f(x) + \epsilon,
\end{equation}

where

\begin{equation}
f(x) \sim \mathcal{GP}\left(m(x), k(x,x')\right).
\end{equation}

Here $m(x)$ is the mean function and $k(x,x')$ is the covariance kernel. The kernel controls the smoothness and characteristic scale of variations in the mass--radius relation. Unlike the piecewise-linear model, the GP does not directly return planetary-regime slopes, but it can reveal whether the data prefer a smooth non-linear relation.

\subsection{Uncertainty Model}

The GP analysis was performed with \texttt{scikit-learn} \cite{pedregosa2011scikit}. The target variable was again

\begin{equation}
y = \log_{10}\left(\frac{R}{R_\oplus}\right),
\end{equation}

as a function of

\begin{equation}
x = \log_{10}\left(\frac{M}{M_\oplus}\right).
\end{equation}

The measured radius uncertainty was included as a heteroscedastic noise term. The mass uncertainty was propagated into the radius direction using the local derivative of the GP mean. Since the derivative is not known before the GP is fitted, the code first initializes the slope using a linear fit and then iteratively refits the GP after updating the propagated mass-error term:

\begin{equation}
\sigma_{\mathrm{obs},i}^2
=
\sigma_{y,i}^2
+
\left[
\frac{d\mu_{\mathrm{GP}}}{dx}(x_i)\sigma_{x,i}
\right]^2.
\end{equation}

In addition, the kernel includes a white-noise term,

\begin{equation}
k(x,x')
=
k_{\mathrm{smooth}}(x,x')
+
\sigma_{\mathrm{white}}^2 \delta_{x,x'},
\end{equation}

which acts as a global intrinsic-scatter parameter. This differs from the Bayesian broken-power-law model, where the intrinsic scatter is allowed to vary by segment.

\subsection{Kernel Comparison}

For a defensible predictive evaluation, the cleaned dataset was split in the same way as the Random Forest analysis: 570 planets were used for training and 190 planets were held back as an untouched final test set. The split was stratified by the three Bayesian mass-regime breakpoints, giving low-mass, intermediate-mass, and giant-planet counts of 69, 203, and 298 in the training set, and 23, 68, and 99 in the test set.

Four covariance kernels were tested: Radial Basis Function (RBF), Matern $3/2$, Matern $5/2$, and Rational Quadratic (RQ). For each kernel, hyperparameters were optimized by maximizing the GP log marginal likelihood on the relevant training fold,

\begin{equation}
\log p(y|X,\theta_{\mathrm{GP}})
=
-\frac{1}{2}(y-m)^T K^{-1}(y-m)
-\frac{1}{2}\log |K|
-\frac{N}{2}\log(2\pi),
\end{equation}

where $K$ includes the smooth kernel, the white-noise term, and the heteroscedastic observational variances. This is an empirical-Bayes procedure: the predictive uncertainties are conditional on the optimized hyperparameters rather than marginalized over a full hyperparameter posterior.

Kernel selection used only the training set. A stratified five-fold cross-validation loop was run on the 570 training planets, and the primary selection metric was the validation negative log predictive density,

\begin{equation}
\mathrm{NLPD}
=
\frac{1}{N_{\mathrm{val}}}
\sum_i
\frac{1}{2}
\left[
\log(2\pi\sigma_i^2)
+
\frac{(y_i-\mu_i)^2}{\sigma_i^2}
\right].
\end{equation}

This score rewards both accurate mean predictions and calibrated predictive variance. RMSE, residual coverage, PTE, and the optimized white-noise level were saved as supporting diagnostics. The final test set was not used during this model-selection step.

\begin{table}[H]
\centering
\caption{Gaussian Process kernel selection using stratified five-fold cross-validation on the training set. Lower NLPD and RMSE are preferred. The final test set is not used in this table.}
\label{tab:gp-model-comparison}
\begin{tabular}{lcccccc}
\toprule
Kernel & CV NLPD & CV RMSE & CV PTE & $f_{1\sigma}$ & $f_{2\sigma}$ & Selected \\
\midrule
RQ & -0.9975 & 0.0904 & 0.352 & 0.695 & 0.960 & yes \\
RBF & -0.9973 & 0.0905 & 0.346 & 0.693 & 0.960 & no \\
Matern $5/2$ & -0.9956 & 0.0906 & 0.352 & 0.693 & 0.960 & no \\
Matern $3/2$ & -0.9907 & 0.0911 & 0.323 & 0.684 & 0.958 & no \\
\bottomrule
\end{tabular}
\end{table}

The Rational Quadratic kernel gives the lowest validation NLPD and is therefore selected for the final GP run. However, this should not be over-interpreted: the difference between RQ and RBF is only $1.5\times10^{-4}$ in mean CV NLPD, much smaller than the pointwise standard error of about $0.04$. All kernels also give similar validation RMSE and two-sigma coverage. The defensible conclusion is therefore not that the data strongly identify the RQ kernel, but that the broad GP result is insensitive to the precise smooth kernel choice.

\subsection{Results}

After kernel selection was fixed, the Rational Quadratic GP was refit on the full 570-object training set and evaluated once on the untouched 190-object test set. The optimized kernel is

\begin{equation}
k(x,x')
=
0.624^2\,k_{\mathrm{RQ}}(\alpha=1.48,\ell=1.72)
+
\sigma_{\mathrm{white}}^2,
\end{equation}

where $\sigma_{\mathrm{white}}=0.083$ dex. The final test metrics are $\mathrm{RMSE}=0.1069$ dex, $\mathrm{MAE}=0.0785$ dex, $R^2=0.917$, and $\mathrm{NLPD}=-0.775$. The residual coverage is 0.668 within one total standard deviation and 0.926 within two total standard deviations.

\begin{figure}[H]
\centering
\includegraphics[width=0.90\linewidth]{gp_mass_radius_best_fit}
\caption{Gaussian Process fit using the selected Rational Quadratic kernel. The training data are shown in grey and the held-out test data in red. The black line is the GP mean and the blue bands show the GP model predictive bands.}
\label{fig:gp-fit}
\end{figure}

The held-out predictive check for the selected GP gives PTE = 0.000, shown in Figure~\ref{fig:gp-ppd}. This means that the observed test-set chi-square statistic is larger than all 1000 replicated statistics in this run. Therefore, although the one-sigma coverage is close to the nominal Gaussian value, the two-sigma coverage and PTE indicate that the held-out predictive distribution is still somewhat too narrow for the most discrepant planets.

\begin{figure}[H]
\centering
\includegraphics[width=0.76\linewidth]{gp_mass_radius_best_ppd_pte}
\caption{Held-out predictive check for the selected Rational Quadratic Gaussian Process model. The PTE value is 0.000.}
\label{fig:gp-ppd}
\end{figure}

\subsection{Kernel-Hyperparameter Sampling Check}

The preceding GP fit is empirical Bayes: the kernel hyperparameters are optimized and then treated as fixed. Since GP predictive bands can be too narrow when hyperparameter uncertainty is ignored, an additional robustness check was implemented using \texttt{emcee}. After the RQ kernel was selected and refit, the propagated measurement-error term was frozen and the kernel hyperparameters were sampled with weak finite log-uniform priors:

\begin{equation}
\sigma_f \in [0.03,2.0], \qquad
\ell \in [0.05,10.0], \qquad
\sigma_{\mathrm{white}} \in [0.005,0.5], \qquad
\alpha_{\mathrm{RQ}} \in [0.05,50.0].
\end{equation}

For each sampled hyperparameter set, the GP log marginal likelihood was computed directly using a Cholesky factorization of

\begin{equation}
K =
K_{\mathrm{RQ}}(\theta)
+
\mathrm{diag}\left(\sigma_{\mathrm{obs},i}^2+\sigma_{\mathrm{white}}^2\right).
\end{equation}

Predictions were then averaged over posterior hyperparameter samples. This marginalizes over amplitude, length scale, white scatter, and the RQ shape parameter while keeping the rest of the analysis unchanged.

\begin{table}[H]
\centering
\caption{Optimized GP and hyperparameter-marginalized GP comparison on the same held-out test set. The marginalized row is a sensitivity check; the current saved diagnostic uses 160 posterior hyperparameter samples and 60 samples for the prediction average.}
\label{tab:gp-hyperparameter-sampling}
\begin{tabular}{lccccccc}
\toprule
Model & RMSE & NLPD & PTE & $f_{1\sigma}$ & $f_{2\sigma}$ & $\sigma_{\mathrm{white}}$ & $\ell$ \\
\midrule
Optimized RQ GP & 0.1069 & -0.775 & 0.000 & 0.668 & 0.926 & 0.083 & 1.72 \\
Hyperparameter-marginalized RQ GP & 0.1069 & -0.770 & 0.000 & 0.668 & 0.926 & 0.082 & 1.59 \\
\bottomrule
\end{tabular}
\end{table}

The marginalization check leaves the point-prediction accuracy and residual coverage almost unchanged, and the PTE remains 0.000. This suggests that the held-out tail under-dispersion is not mainly caused by fixing the RQ kernel hyperparameters. A more likely limitation is the simple global Gaussian white-scatter model, or the fact that the GP uses only mass as a predictor even though radius also depends on composition, irradiation, age, and observational selection.

\subsection{Discussion}

The GP model provides a useful non-parametric comparison to the Bayesian broken-power-law model. It captures the global curvature of the mass--radius relation without specifying the number or position of planetary regimes, and the stratified train-test split gives a more honest estimate of predictive performance than a full-sample residual calculation. The selected GP is competitive as a smooth probabilistic benchmark, but its held-out predictive check shows that the uncertainty model is not perfect.

The main disadvantage is interpretability. The GP does not directly identify physically meaningful slopes for terrestrial, intermediate, and giant planets. Its extrapolation behaviour is also less physically motivated: away from the data, the GP prediction tends to revert toward the learned mean structure according to the chosen kernel. The hyperparameter-sampling check makes the uncertainty treatment more defensible, but it does not remove the remaining tail-calibration issue. Therefore, the GP is valuable as a flexible probabilistic regression tool, while the broken-power-law model remains more directly connected to planetary-structure interpretation.

\section{Random Forest Regression}

\subsection{Motivation}

Random Forests are ensemble machine-learning models based on many decision trees trained on bootstrap samples of the data \cite{breiman2001random}. They are attractive as a predictive baseline because they can learn non-linear structure without assuming a parametric functional form. In this project, the Random Forest is used to fit

\begin{equation}
\log_{10}\left(\frac{R}{R_\oplus}\right)
=
f\left[
\log_{10}\left(\frac{M}{M_\oplus}\right)
\right].
\end{equation}

Unlike the Bayesian and GP methods, the Random Forest does not naturally provide a probabilistic posterior distribution. Therefore, the standard deviation of predictions across individual trees was used as a practical uncertainty proxy. This should be interpreted cautiously: tree-to-tree scatter is an ensemble-dispersion measure, not a full Bayesian predictive uncertainty.

\subsection{Training and Hyperparameter Selection}

The cleaned dataset was split into a training set and an untouched held-out test set with test fraction 0.25. The split was stratified by the three broad mass regimes identified by the baseline Bayesian model, using breakpoints at $\log_{10}M=0.824$ and $\log_{10}M=2.196$. This gives 570 training planets and 190 test planets. The corresponding low-mass, intermediate-mass, and giant-planet counts are 69, 203, and 298 in the training set, and 23, 68, and 99 in the test set. Stratifying the split avoids evaluating the model on a test set that is accidentally dominated by one planetary regime.

The training used inverse-variance sample weights based on the measured log-radius uncertainty, with a small error floor to avoid giving extreme weight to a few very precise measurements. Several smoothness settings were compared. The main hyperparameters were the tree depth, the minimum number of samples per leaf, and the bootstrap fraction used for each tree. Hyperparameters were selected using only the training set through out-of-bag (OOB) predictions. For each training planet, the OOB prediction is formed from trees whose bootstrap sample did not include that planet. This provides an internal validation estimate without spending additional data on a separate validation set.

The OOB objective was RMSE in log-radius space:

\begin{equation}
\mathrm{RMSE}_{\mathrm{OOB}}
=
\sqrt{
\frac{1}{N_{\mathrm{train}}}
\sum_i
\left(y_i-\hat{y}_{i,\mathrm{OOB}}\right)^2
}.
\end{equation}

The final test set was not used during this hyperparameter scan. The model with the absolute best OOB RMSE was not automatically adopted, because deeper forests can produce visually jagged fits and overfit local fluctuations. Instead, the final run used a one-standard-error selection rule: first identify the best OOB RMSE, then choose the smoothest model whose OOB RMSE is statistically consistent with that best value.

\begin{table}[H]
\centering
\caption{Selected Random Forest hyperparameter comparison using training-set out-of-bag predictions. The final model is chosen by the one-standard-error rule, not by the absolute minimum OOB RMSE.}
\label{tab:rf-model-comparison}
\begin{tabular}{lcccccc}
\toprule
Model & Depth & Min. leaf & Max samples & OOB RMSE & OOB SE & Selected \\
\midrule
medium & 5 & 12 & 0.8 & 0.0950 & 0.0037 & no \\
medium full & 5 & 12 & 1.0 & 0.0953 & 0.0037 & no \\
balanced & 6 & 8 & 0.8 & 0.0956 & 0.0036 & no \\
smooth leaf 18 full & 4 & 18 & 1.0 & 0.0961 & 0.0037 & no \\
smooth full & 4 & 24 & 1.0 & 0.0984 & 0.0037 & yes \\
ultra smooth & 2 & 40 & 0.6 & 0.1210 & 0.0049 & no \\
\bottomrule
\end{tabular}
\end{table}

The absolute best OOB RMSE was obtained by the \texttt{medium} model, with OOB RMSE = 0.0950 dex. The one-standard-error threshold was 0.0986 dex, and the selected smoother model was \texttt{smooth\_full}, with max depth 4, minimum leaf size 24, and full bootstrap samples. Its OOB RMSE was 0.0984 dex. After this selection was fixed, the selected model was evaluated once on the untouched test set.

\subsection{Uncertainty Proxy and Predictive Checks}

For a fitted forest with $T$ trees, the mean prediction is

\begin{equation}
\hat{y}(x)
=
\frac{1}{T}
\sum_{t=1}^{T}
\hat{y}_t(x),
\end{equation}

and the tree-spread proxy is

\begin{equation}
\sigma_{\mathrm{tree}}^2(x)
=
\frac{1}{T-1}
\sum_{t=1}^{T}
\left[\hat{y}_t(x)-\hat{y}(x)\right]^2.
\end{equation}

The total variance used for residual coverage and the approximate predictive check was

\begin{equation}
\sigma_{\mathrm{tot},i}^2
=
\sigma_{\mathrm{tree},i}^2
+
\sigma_{y,i}^2
+
\sigma_{\mathrm{prop},i}^2,
\end{equation}

where $\sigma_{\mathrm{prop},i}$ is the finite-difference propagation of the measured mass uncertainty through the Random Forest mean prediction.

Figure~\ref{fig:rf-fit} shows the selected Random Forest fit. The training data are shown in grey and the held-out test data in red. The final curve is smooth relative to the more flexible forests, but it still has the piecewise-constant character typical of tree-based regression.

\begin{figure}[H]
\centering
\includegraphics[width=0.90\linewidth]{rf_mass_radius_best_fit}
\caption{Random Forest fit for the selected \texttt{smooth\_full} model. The black curve is the mean prediction across trees and the blue bands show the raw tree-spread uncertainty proxy.}
\label{fig:rf-fit}
\end{figure}

The selected model gives train RMSE = 0.0911 dex, OOB RMSE = 0.0984 dex, and test RMSE = 0.1192 dex, with test $R^2=0.896$. This indicates useful but not leading predictive performance on held-out data. The clean test-set estimate is less optimistic than the earlier exploratory comparison because the test set was no longer used during model selection. The uncertainty calibration is still poor: the predictive check gives PTE = 0.000, and only 0.353 of test planets lie within one total standard deviation while 0.595 lie within two total standard deviations.

\begin{figure}[H]
\centering
\includegraphics[width=0.76\linewidth]{rf_mass_radius_best_ppd_pte}
\caption{Approximate predictive check for the selected Random Forest model. The PTE value is 0.000, showing that the raw tree-spread uncertainty proxy is too narrow to reproduce the observed test-set scatter.}
\label{fig:rf-ppd}
\end{figure}

\subsection{Discussion}

The Random Forest is useful as a predictive machine-learning baseline. It tests how well a flexible empirical method can predict radii without imposing a parametric mass--radius shape. The stratified train-test split makes overfitting visible, while OOB model selection keeps the final test set reserved for a single unbiased evaluation. The hyperparameter scan also shows the expected bias-variance tradeoff: very shallow forests underfit, while deeper forests reduce training error but can produce a less physically plausible, locally jagged relation.

The main weakness is uncertainty quantification. The tree-spread band is an ensemble-disagreement proxy, not a calibrated predictive interval for this problem. Individual trees in a Random Forest are correlated because they are trained on overlapping bootstrap samples and the same one-dimensional predictor. As a result, the tree-to-tree dispersion substantially underestimates the astrophysical scatter in radius at fixed mass. Therefore, Random Forests are useful for empirical point prediction, while the Bayesian broken-power-law model and GP provide more explicit probabilistic alternatives. A natural future improvement would be to calibrate the Random Forest intervals using conformal prediction or a validation-based interval calibration.

\section{Shared Held-Out Predictive Comparison with \texttt{compare\_MR\_fit\_performance.py}}

The script \texttt{compare\_MR\_fit\_performance.py} compares the Bayesian broken-power-law model, Gaussian Process, and Random Forest on exactly the same mass-regime-stratified held-out test set. The cleaned DACE sample was split into 570 training planets and 190 test planets using the same broad Bayesian mass regimes as the GP and RF sections. The split identifiers are saved so every method is evaluated on the same planets. The final test set was not used to choose the number of Bayesian breakpoints, the GP kernel, the RF configuration, or any hyperparameters. The two-breakpoint \texttt{emcee} model and \texttt{smooth\_full} forest were pre-declared baselines; all their fitted parameters were learned from the training data only. The GP kernel was selected by stratified cross-validation within the training set using NLPD.

The comparison uses RMSE to measure point-prediction accuracy in log-radius dex. It also reports Gaussian negative log predictive density (NLPD) in log-radius space using each method's predictive mean and variance for the observed held-out log-radius values. For the \texttt{emcee} model, the posterior predictive mixture was moment-matched to a Gaussian, so its NLPD is a common-score approximation rather than the exact mixture log score. For the GP, the existing predictive-variance helper was used so that the fitted \texttt{WhiteKernel} contribution was not added twice. For the RF, NLPD is only a diagnostic: the variance uses raw tree-to-tree spread, radius measurement error, and the existing finite-difference mass-error propagation proxy, but this is not a calibrated posterior predictive distribution.

\begin{table}[H]
\centering
\caption{Shared held-out comparison for the two-breakpoint \texttt{emcee} model, selected-kernel GP, and selected Random Forest. RMSE and MAE are in dex in $\log_{10}R$. NLPD is a Gaussian score in log-radius space; the RF value is proxy-only. Values are taken from \texttt{mr\_methods\_shared\_split\_comparison.csv}.}
\label{tab:shared-heldout-comparison}
\begin{tabular}{lccccccc}
\toprule
Method & RMSE & Radius frac. & NLPD & MAE & $R^2$ & $f_{1\sigma}$ & $f_{2\sigma}$ \\
\midrule
\texttt{emcee\_2break} & 0.1063 & 0.277 & -0.864 & 0.0785 & 0.918 & 0.658 & 0.921 \\
\texttt{gp\_selected\_kernel} & 0.1069 & 0.279 & -0.775 & 0.0785 & 0.917 & 0.668 & 0.926 \\
\texttt{rf\_smooth\_full} proxy & 0.1192 & 0.316 & 3.597 & 0.0866 & 0.896 & 0.353 & 0.595 \\
\bottomrule
\end{tabular}
\end{table}

\begin{figure}[H]
\centering
\includegraphics[width=0.88\linewidth]{mr_methods_shared_split_rmse_nlpd}
\caption{Shared held-out RMSE and Gaussian NLPD comparison. The Random Forest NLPD is marked as a proxy because it uses raw tree-to-tree scatter rather than a calibrated predictive distribution.}
\label{fig:shared-heldout-rmse-nlpd}
\end{figure}

This shared comparison gives a clean summary of the tradeoff. The \texttt{emcee} and GP models have nearly identical held-out RMSE, while the two-breakpoint \texttt{emcee} baseline has the better Gaussian NLPD in this run. The RF is still useful as a non-parametric point-prediction benchmark, but its raw tree-spread uncertainty is much too narrow: both the high proxy NLPD and the low one- and two-sigma coverage show that it should not be interpreted as a calibrated uncertainty model. RMSE is therefore the fairest cross-method point-prediction comparison, while NLPD is most meaningful for the Bayesian and GP predictive distributions.

\section{Conclusion}

This work compared four statistical approaches to the exoplanet mass--radius relation using the same cleaned DACE sample of 760 planets. All methods model $\log_{10}R=f(\log_{10}M)$. This makes the Bayesian slopes directly interpretable as power-law exponents and gives the machine-learning methods a stable one-dimensional regression problem.

The broken-power-law model provides the clearest physical interpretation. In the baseline two-breakpoint MCMC fit, the inferred breakpoints are near $6.7\,M_\oplus$ and $152\,M_\oplus$, separating small, intermediate, and giant planets. The slopes are approximately $0.40$, $0.61$, and $-0.04$: a positive small-planet scaling, rapid radius growth in the intermediate regime, and weak mass dependence for gas giants. Segment-dependent intrinsic scatter is essential because measurement errors alone cannot explain the observed spread at fixed mass.

The preferred breakpoint count depends on the statistical question. BIC favours two breakpoints through a strong complexity penalty; WAIC gives more support to three or four breakpoints for predictive fit; and the current nested-sampling evidence run favours two breakpoints after integrating over each model's prior volume. These criteria are not interchangeable. Together they support multiple planetary regimes while retaining the two-breakpoint model as the most parsimonious physical description.

The GP is a strong non-parametric probabilistic alternative. It follows the global curvature without explicit breakpoints and, on the shared held-out set, its RMSE of 0.1069 dex is nearly identical to the two-breakpoint model's 0.1063 dex. Its PTE nevertheless shows tail under-dispersion that is not resolved by marginalizing the selected kernel's hyperparameters. The Random Forest provides a useful empirical benchmark, but its held-out RMSE is higher at 0.1192 dex and its raw tree-spread proxy strongly underestimates predictive scatter.

Overall, the \texttt{emcee} broken-power-law analysis is the strongest choice when physical interpretation and calibrated uncertainty are central. The \texttt{dynesty} evidence analysis complements it by comparing complete Bayesian models rather than posterior point estimates. The GP is the strongest flexible alternative for smooth probabilistic prediction, while the Random Forest should be treated as a point-prediction benchmark unless its intervals are independently calibrated. The shared evaluation is crucial: it separates these predictive conclusions from the full-sample diagnostics used to interpret the planetary regimes.

\printbibliography

\end{document}
